\begin.Abstract}
We present the Bailout Protocol, a foundational safety architecture for multi-agent systems operating under the BX3 Framework. The protocol defines a strict boundary condition: whenever any agent node encounters a state it cannot resolve within its operational boundaries, it must escalate accountability upward through the hierarchy rather than defaulting to autonomous resolution. This escalation continues until it reaches a designated human anchor, who bears ultimate accountability for the outcome. The core thesis of this work holds that \textbf{the system fails upward into human consciousness. It never fails downward into algorithmic chaos.} We argue that unconstrained algorithmic failure cascades represent an existential risk in autonomous deployments, and that human escalation is not a weakness of hybrid systems but rather their primary safety mechanism. The Bailout Protocol formalizes this escalation as a first-class protocol with defined triggering conditions, routing logic, timeout behavior, and audit requirements. We present a reference implementation within AgentOS, demonstrate its behavior across cascading boundary conditions, and analyze its implications for accountability, governance, and trust in agentic deployments. The protocol contributes a concrete operational layer to the broader BX3 Framework, distinguishing it from prior human-in-the-loop approaches by embedding accountability escalation as a mandatory architectural constraint rather than an optional safeguard or fallback heuristic. Results indicate that systems employing the Bailout Protocol maintain human coherence under uncertainty significantly longer than fully autonomous baselines.
\end.Abstract}